% ---
% id: receiver/component-catalogue/def:c10-tolerance
% kind: def
% title: "Admission test of the generic block"
% summary: >-
%   The contractivity of the generic block C10 is decided by the route-(ii)
%   test of [[def:tolmap]] (v), at the same widened threshold, and its lossless
%   case, an equality, is decided two-sidedly within ε₁.
%   Without the widening, rounding one unit above one would refuse a unitary
%   block.
% prominence: standard
% status: accepted
% uses: [receiver/def:tolmap, thm:catalogue-admissibility]
% mentions: [ex:bb84, def:catalogue-entries]
% assumes: []
% refs: []
% provenance: newly-derived
% lean:
% ---

The generic block is the one catalogue entry whose admission is a spectral
test on a supplied matrix, and that test is evaluated in floating point. It
is decided at $\varepsilon_1$, one of the \tref{receiver tolerances}, for the
same reason as at route (ii): a lossless block has
$\lmax(\Gmat_\Sigma^\dagger\Gmat_\Sigma)=1$ \emph{exactly} (the worked
instrument \eqref{eq:bb84G} of Example~\ref{ex:bb84} is unitary), so a
computed value one rounding unit above $1$ would refuse it.

\begin{definition}
\label{def:c10-tolerance}
The (E1) predicate of catalogue entry C10 of
Definition~\ref{def:catalogue-entries} is decided by the route-(ii)
contractivity test of Definition~\ref{def:tolmap}\,(v) applied to the supplied
block $\Gmat_\Sigma$: the block is admitted when
$\lmax(\Gmat_\Sigma^\dagger\Gmat_\Sigma)\le1+\varepsilon_1$. C10 is the one
entry whose (E1) is a spectral condition rather than a parameter-domain
comparison (Theorem~\ref{thm:catalogue-admissibility}\,(iii)). Its lossless
sub-case $\Gmat_\Sigma^\dagger\Gmat_\Sigma=\id$ is an \emph{equality} and is
decided two-sidedly, as
$\lVert\Gmat_\Sigma^\dagger\Gmat_\Sigma-\id\rVert\le\varepsilon_1$.
\end{definition}
